% Standalone resolution manuscript generated from solution.md.
% Compile with XeLaTeX; author and review metadata are maintained in Markdown.
\documentclass[10pt,a4paper]{article}
\usepackage[margin=24mm,headheight=15pt]{geometry}
\usepackage{fontspec}
\setmainfont{DejaVuSerif.ttf}[BoldFont=DejaVuSerif-Bold.ttf,ItalicFont=DejaVuSerif-Italic.ttf,BoldItalicFont=DejaVuSerif-BoldItalic.ttf]
\setsansfont{DejaVuSans.ttf}[BoldFont=DejaVuSans-Bold.ttf,ItalicFont=DejaVuSans-Oblique.ttf]
\setmonofont{DejaVuSansMono.ttf}
\usepackage{amsmath,amssymb,unicode-math}
\setmathfont{latinmodern-math.otf}
\usepackage{xcolor,microtype,parskip,needspace,fancyhdr}
\usepackage{bookmark}
\usepackage{xurl}
\hypersetup{colorlinks=true,urlcolor=blue,linkcolor=blue,pdftitle={MF-18:
Imaginary-part rank for the general complex Green-function
equation},pdfauthor={George Stepaniants},pdflang={en-GB}}
\urlstyle{same}
\setlength{\emergencystretch}{3em}
\providecommand{\tightlist}{\setlength{\itemsep}{0pt}\setlength{\parskip}{0pt}}
\providecommand{\pandocbounded}[1]{#1}
\setcounter{secnumdepth}{0}
\allowdisplaybreaks[1]
\widowpenalty=10000
\clubpenalty=10000
\pagestyle{fancy}
\fancyhf{}
\fancyhead[L]{\small\sffamily RESOLUTION / MF-18}
\fancyhead[R]{\small\sffamily 11 September 2026}
\fancyfoot[L]{\parbox[b]{0.9\textwidth}{\fontsize{8}{10}\selectfont AI-assisted
proof; independent automated-agent review documented in the submission
record.}}
\fancyfoot[R]{\footnotesize\thepage}
\begin{document}
{\raggedright\Large\bfseries MF-18: Imaginary-part rank for the general
complex Green-function equation\par}
\vspace{0.4em}
{\normalsize\bfseries George Stepaniants\par}
{\small Department of Computing and Mathematical Sciences, California
Institute of Technology, Pasadena, California, USA.\par}
\vspace{0.6em}
This manuscript proves the complete complex-coefficient rank equality in
MF-18. It was developed with substantial ChatGPT/Codex assistance at the
author's request. The accompanying independent review is automated-agent
review, not external human peer review or formal proof-assistant
verification.

\subsection{Theorem 1}\label{theorem-1}

Let \(n\geq1\) and let \(C,D,R,P\in\mathbb C^{n\times n}\), with
\(R=R^*\), \(P=P^*\), and

\[
P+\lambda D^*+\lambda^{-1}D\succ0\qquad(|\lambda|=1).
\]

For every \(\eta>0\), suppose \(X_\eta\) is a nonsingular solution of

\[
X_\eta+(C^*+i\eta D^*)X_\eta^{-1}(C+i\eta D)=R+i\eta P
\]

such that \(\rho(X_\eta^{-1}(C+i\eta D))<1\). Suppose \(X_\eta\) tends
to a finite nonsingular \(X_0\) as \(\eta\downarrow0\). Suppose the
matrix polynomial

\[
\mathcal P_0(\lambda)=\lambda^2C^*-\lambda R+C
\]

is regular and all its unit-circle eigenvalues are algebraically simple.
If their number is \(2m\), then

\[
\boxed{\operatorname{rank}\left(\frac{X_0-X_0^*}{2i}\right)=m.}
\]

The proof permits general complex \(C,D\), Hermitian \(R,P\), singular
\(C\) or \(D\), unit-circle roots at \(1\) or \(-1\), and arbitrary
Jordan structure strictly inside the unit circle. The uniqueness assumed
in the canonical statement is not needed beyond specifying the given
stabilizing family.

\subsection{1. Selecting the stable polynomial
eigenvalues}\label{selecting-the-stable-polynomial-eigenvalues}

Write

\[
A_\eta=C+i\eta D,\qquad B_\eta=C^*+i\eta D^*,\qquad
Q_\eta=R+i\eta P,\qquad S_\eta=X_\eta^{-1}A_\eta,
\]

and

\[
\mathcal P_\eta(\lambda)=\lambda^2B_\eta-\lambda Q_\eta+A_\eta.
\]

The finite nonsingular limit implies \(S_\eta\to S=X_0^{-1}C\), so
\(\rho(S)\leq1\) by continuity of the roots of a characteristic
polynomial. Notice that \(B_\eta\) is generally different from
\(A_\eta^*\).

For \(|\lambda|=1\), the positivity assumption at \(-\lambda\) gives

\[
W(\lambda):=P-\lambda D^*-\lambda^{-1}D\succ0. \tag{1}
\]

Also \(P\succ0\), since it is the average of the assumed positive
matrices at \(\lambda\) and \(-\lambda\).

\textbf{Lemma 2.} For each \(\eta>0\), the scalar polynomial
\(\det\mathcal P_\eta\) has exactly \(n\) zeros in the open unit disk,
counted with algebraic multiplicity, and none on its boundary.

\textbf{Proof.} For \(t\in[0,1]\), consider

\[
\mathcal P_{\eta,t}(\lambda)
=\lambda^2(tC^*+i\eta tD^*)
 -\lambda(tR+i\eta P)+(tC+i\eta tD).
\]

On the unit circle,

\[
\lambda^{-1}\mathcal P_{\eta,t}(\lambda)
=t(\lambda C^*+\lambda^{-1}C-R)
 -i\eta\bigl[P-t(\lambda D^*+\lambda^{-1}D)\bigr].
\]

The first term is Hermitian, and the bracket is the positive-definite
matrix \((1-t)P+tW(\lambda)\). The displayed matrix has
negative-definite Hermitian imaginary part, hence is nonsingular. Its
determinant therefore has no zero on the unit circle throughout the
homotopy. The argument principle shows that the number of zeros inside
is constant in \(t\). At \(t=0\), the polynomial is \(-i\eta\lambda P\),
whose determinant has exactly \(n\) zeros at zero. This proves the lemma
without requiring either coefficient of \(\mathcal P_\eta\) to be
invertible. \(\square\)

The nonlinear equation gives the exact factorization

\[
\mathcal P_\eta(\lambda)
=(\lambda B_\eta X_\eta^{-1}-I)X_\eta(\lambda I-S_\eta). \tag{2}
\]

Its coefficients are \(B_\eta\), \(-(B_\eta S_\eta+X_\eta)=-Q_\eta\),
and \(X_\eta S_\eta=A_\eta\). The last determinant factor already
contributes \(n\) zeros in the open unit disk, because \(S_\eta\) is
stabilizing. By Lemma 2, the polynomial

\[
\ell_\eta(\lambda)=\det(\lambda B_\eta X_\eta^{-1}-I)
\]

has no zero in that disk.

Its coefficients converge to those of

\[
\ell_0(\lambda)=\det(\lambda C^*X_0^{-1}-I),
\]

and \(\ell_0(0)=(-1)^n\ne0\). Thus \(\ell_0\) is not identically zero
and also has no zero in the open unit disk. Indeed, if it had such a
zero, choose a small circle surrounding that zero, contained in the
disk, and containing no other zero on its boundary. Uniform convergence
on that circle and Rouché's theorem would force \(\ell_\eta\) to have a
zero inside for all sufficiently small \(\eta\), a contradiction.

Taking limits in (2) now gives

\[
\mathcal P_0(\lambda)
=(\lambda C^*X_0^{-1}-I)X_0(\lambda I-S). \tag{3}
\]

Consequently every eigenvalue of \(\mathcal P_0\) strictly inside the
unit disk is an eigenvalue of \(S\), with exactly the same algebraic
multiplicity there.

\subsection{2. Counting the eigenvalues at the
limit}\label{counting-the-eigenvalues-at-the-limit}

Let \(p(\lambda)=\det\mathcal P_0(\lambda)\). The coefficient symmetry
gives, for \(\lambda\ne0\),

\[
\mathcal P_0(\lambda)
=\lambda^2\mathcal P_0(1/\overline\lambda)^*,
\qquad
p(\lambda)=\lambda^{2n}\overline{p(1/\overline\lambda)}. \tag{4}
\]

Write \(d=\deg p\) and let \(z\) be its vanishing order at zero. The
polynomial is nonzero by regularity. Comparing the lowest and highest
nonzero coefficients in (4) gives

\[
d+z=2n. \tag{5}
\]

Each nonzero root off the unit circle is paired, with the same
multiplicity, with its reciprocal conjugate. There are \(2m\)
unit-circle roots, so exactly half of the remaining \(d-z-2m\) nonzero
roots are strictly inside the disk. Including the zero root, the number
of roots strictly inside is therefore

\[
z+\frac{d-z-2m}{2}=n-m. \tag{6}
\]

This argument accounts for possible singular leading coefficients: the
degree \(d\) need not equal \(2n\).

By (3), \(S\) has exactly \(n-m\) eigenvalues strictly inside the unit
disk, counted with multiplicity. Since its order is \(n\) and
\(\rho(S)\leq1\), its remaining \(m\) eigenvalues lie on the unit
circle. They are all distinct and algebraically simple because every
corresponding zero of \(p\) is simple. In particular, the unit-circle
part of \(S\) is semisimple, and the powers \(S^j\) are bounded for
\(j\geq0\).

\subsection{3. A Stein identity and the missing lower
bound}\label{a-stein-identity-and-the-missing-lower-bound}

Set \(X=X_0\) and \(H=(X-X^*)/(2i)\). Taking the limit of the nonlinear
equation gives

\[
X+C^*X^{-1}C=R.
\]

Since

\[
\frac{X^{-1}-X^{-*}}{2i}=-X^{-*}HX^{-1},
\]

taking Hermitian imaginary parts yields the exact Stein identity

\[
H=S^*HS. \tag{7}
\]

Let \(\lambda=e^{i\theta_0}\) be one of the \(m\) unit-circle
eigenvalues of \(S\), and let \(v\) be a corresponding unit eigenvector.
Then

\[
Cv=\lambda Xv. \tag{8}
\]

Consider the Hermitian matrix-valued function

\[
F(\theta)=e^{i\theta}C^*+e^{-i\theta}C-R.
\]

We have \(\mathcal P_0(e^{i\theta})=e^{i\theta}F(\theta)\). Thus the
simple determinant root of \(\mathcal P_0\) at \(\lambda\) gives a
simple zero of \(\det F(\theta)\) at \(\theta_0\). It follows that
\(F(\theta_0)\) has one-dimensional kernel spanned by \(v\), and

\[
a:=v^*F'(\theta_0)v
=v^*(i\lambda C^*-i\lambda^{-1}C)v\ne0. \tag{9}
\]

For completeness, a simple determinant zero forces nullity one because a
matrix with nullity at least two has zero adjugate and therefore zero
determinant derivative. For the Hermitian matrix \(F(\theta_0)\) with
one-dimensional kernel, its adjugate is \(\kappa vv^*\) for a nonzero
real \(\kappa\). The determinant derivative is consequently
\(\kappa v^*F'(\theta_0)v\), proving (9), including when \(n=1\).

Equation (8) and \(|\lambda|=1\) give the direct identity

\[
a=i\bigl(v^*X^*v-v^*Xv\bigr)=2v^*Hv. \tag{10}
\]

Hence \(v^*Hv\ne0\) at each of the selected unit-circle eigenvectors. If
\(v_j,v_k\) correspond to distinct unit-circle eigenvalues
\(\lambda_j,\lambda_k\), (7) gives

\[
(1-\overline{\lambda_j}\lambda_k)v_j^*Hv_k=0.
\]

The scalar prefactor is nonzero for \(j\ne k\), so all these cross terms
vanish. The Gram matrix of \(H\) on the span of the \(m\) unit-circle
eigenvectors is diagonal with nonzero diagonal entries, by (10). It is
nonsingular, and therefore

\[
\operatorname{rank}H\geq m. \tag{11}
\]

\subsection{4. The stable generalized eigenspace gives the upper
bound}\label{the-stable-generalized-eigenspace-gives-the-upper-bound}

Let \(E_s\) be the generalized eigenspace of \(S\) corresponding to its
eigenvalues strictly inside the unit disk. Its dimension is \(n-m\), and
\(S^j x\to0\) for every \(x\in E_s\). Iterating (7) gives, for any
\(y\in\mathbb C^n\),

\[
x^*Hy=(S^j x)^*H S^j y.
\]

The right-hand side tends to zero because \(S^j x\to0\) and the powers
of \(S\) are bounded. Hence \(x^*H=0\). Since \(H\) is Hermitian,
\(Hx=0\), and \(E_s\subseteq\ker H\). Consequently

\[
\operatorname{rank}H\leq m. \tag{12}
\]

Together, (11) and (12) prove Theorem 1. If \(m=0\), the stable-subspace
argument directly gives \(H=0\). No differentiability of the family
\(X_\eta\), positivity of \(H\), or diagonalizability of the stable part
was used. \(\square\)

\subsection{Scope, attribution, and
status}\label{scope-attribution-and-status}

Guo, Kuo, and Lin prove the upper bound and formulate the equality
conjecture in their 2012 JCAM paper. Their earlier SIMAX paper addresses
the real-coefficient scalar-regularization setting. Matthew J.
Colbrook's existing repository contribution proves an auxiliary
real-coefficient theorem that also treats defective unit-circle
structure. That contribution remains credited and is not a premise of
the present general-complex proof. The argument here retains the
canonical simple-eigenvalue hypothesis and does not claim the defective
extension.

The public-network audit at 22:37 UTC on 11 September 2026 checked five
repositories and all 29 public branch heads. Every checked canonical
MF-18 page retained partial status; the matching discussions concerned
the existing real-coefficient auxiliary result. A bounded primary-source
literature check likewise found the conjecture, known upper bound, and
prior real subcase. The
\href{https://github.com/ajt60gaibb/OpenProblemsInNLA/blob/main/references/stepaniants-mf18-2026-09-11/README.md}{submission
record} documents the audit, mathematical review, and exact reviewed
source. These checks do not certify historical priority.

\subsection{References}\label{references}

\begin{enumerate}
\def\labelenumi{\arabic{enumi}.}
\tightlist
\item
  \href{https://github.com/ajt60gaibb/OpenProblemsInNLA/blob/aaa88c4/matrix-functions-and-stability/MF-18/README.md}{MF-18
  canonical problem statement}.
\item
  C.-H. Guo, Y.-C. Kuo, W.-W. Lin,
  \href{https://doi.org/10.1016/j.cam.2012.05.012}{\emph{Numerical
  solution of nonlinear matrix equations arising from Green's function
  calculations in nano research}}, Journal of Computational and Applied
  Mathematics 236 (2012), 4166--4180; Theorem 5 and the conjecture on
  p.~4172. \href{https://uregina.ca/~chguo/JCAM_GuoKuoLin.pdf}{Author
  manuscript}.
\item
  C.-H. Guo, Y.-C. Kuo, W.-W. Lin,
  \href{https://doi.org/10.1137/100814706}{\emph{On a nonlinear matrix
  equation arising in nano research}}, SIAM Journal on Matrix Analysis
  and Applications 33 (2012), 235--262, Section 3.
  \href{https://uregina.ca/~chguo/simax81470.pdf}{Author manuscript}.
\item
  Matthew J. Colbrook,
  \href{https://github.com/ajt60gaibb/OpenProblemsInNLA/blob/aaa88c4/references/colbrook-matrix-functions-2026-09-11/manuscripts/MF-18.tex}{existing
  MF-18 auxiliary manuscript}, Theorem 1, Corollary 3, and Section 8,
  with the repository's stated real-coefficient scope and
  independent-agent review.
\end{enumerate}
\end{document}
